\chapter{Diverticular Disease Surgery}

\section*{What Is This Surgery?}
Diverticular disease surgery, most commonly a sigmoid colectomy, is a surgical intervention aimed at resecting the segment of the colon affected by diverticula and associated inflammation or complications. This procedure is primarily indicated for patients suffering from complicated diverticulitis, which encompasses conditions such as abscess formation, free perforation leading to peritonitis, fistula development, or bowel obstruction. It is also considered for individuals experiencing recurrent episodes of uncomplicated diverticulitis that significantly impair their quality of life, or for diverticular bleeding refractory to less invasive treatments. The overarching goal of surgery is to eliminate the diseased colonic segment, resolve acute complications, prevent future recurrences, and restore normal bowel function, often through primary anastomosis or, in acute settings, via a temporary or permanent stoma.

\section*{Alternative Names}
\begin{itemize}
  \item Sigmoid Colectomy
  \item Resection for Diverticulitis
  \item Diverticular Resection
  \item Hartmann's Procedure (specifically for emergent cases with significant inflammation or contamination)
  \item Laparoscopic Colectomy for Diverticulitis
  \item Open Colectomy for Diverticulitis
  \item Left Hemicolectomy (if disease extends proximally)
\end{itemize}

\section*{Relevant Surgical Anatomy}
The primary anatomical focus for diverticular disease surgery is the sigmoid colon, which is the most common site for diverticula formation due to its smaller lumen diameter and higher intraluminal pressure. The sigmoid colon is an S-shaped segment of the large intestine, typically 25-40 cm long, extending from the descending colon at the pelvic brim to the rectum at the level of S3. It is characterized by its mesentery, the sigmoid mesocolon, which allows for significant mobility but also contains its vascular supply and lymphatic drainage.

Key anatomical structures and relations include:
\begin{itemize}
  \item \textbf{Vascular Supply:} The sigmoid colon is primarily supplied by the sigmoid arteries, branches of the inferior mesenteric artery (IMA). The IMA also gives rise to the left colic artery (supplying the descending colon) and the superior rectal artery (supplying the upper rectum). The marginal artery of Drummond, a continuous arterial arcade along the mesenteric border of the colon, connects the superior and inferior mesenteric arterial systems and is crucial for maintaining collateral flow, especially during IMA ligation.
  \item \textbf{Venous Drainage:} Venous drainage mirrors the arterial supply, with sigmoid veins draining into the inferior mesenteric vein (IMV), which typically empties into the splenic vein.
  \item \textbf{Lymphatic Drainage:} Lymphatics follow the arterial supply, draining into lymph nodes along the sigmoid arteries, IMA, and ultimately to para-aortic nodes.
  \item \textbf{Nerve Relations:} Autonomic nerves, including the superior and inferior hypogastric plexuses, are located in the retroperitoneum and pelvis. Injury to these nerves during extensive pelvic dissection can lead to bladder or sexual dysfunction.
  \item \textbf{Adjacent Structures:} The left ureter is a critical structure that lies retroperitoneally and crosses the common iliac artery, often running close to the root of the sigmoid mesentery. It must be meticulously identified and protected during mobilization of the left colon. Other adjacent structures include the bladder anteriorly, the small bowel loops, and in females, the uterus and left ovary.
\item \textbf{Surgical Landmarks:} The sacral promontory marks the transition from the descending colon to the sigmoid colon. The peritoneal reflections around the sigmoid mesocolon are important for initiating mobilization.
\end{itemize}

\section*{Surgical Principle \& Rationale}
The fundamental surgical principle for diverticular disease is the complete excision of the segment of colon harboring the diverticula and responsible for the patient's symptoms or complications. This almost invariably involves the sigmoid colon, as it is the most common and often the sole site of clinically significant diverticular disease. By removing the diseased tissue, the surgeon eliminates the source of ongoing infection, inflammation, perforation, bleeding, or obstruction, thereby resolving the acute problem and preventing future episodes.

The rationale extends beyond mere removal; it encompasses restoring bowel continuity and function while minimizing morbidity. Technical goals include achieving a tension-free, well-perfused anastomosis between the remaining healthy bowel segments, ensuring adequate margins of resection to include all diseased diverticula, and meticulously controlling intra-abdominal contamination, especially in emergent settings. In cases of severe inflammation, generalized peritonitis, or patient instability, a primary anastomosis may be deemed unsafe due to an unacceptably high risk of anastomotic leak. In such scenarios, a Hartmann's procedure is often performed, involving resection of the diseased colon, closure of the distal rectal stump, and creation of a diverting colostomy. This allows for resolution of the acute process, with a planned reversal of the colostomy at a later date when the patient is stable and inflammation has subsided.

\section*{History \& Surgical Evolution}
The surgical management of diverticular disease has undergone significant evolution since its early descriptions. Diverticulitis was first recognized as a distinct clinical entity in the late 19th and early 20th centuries. Early surgical approaches were often multi-stage procedures, reflecting the high morbidity and mortality associated with direct resection and primary anastomosis in the presence of acute inflammation and contamination. Surgeons like Mikulicz and Paul developed techniques involving exteriorization of the diseased bowel, followed by delayed resection and stoma closure, to manage the acute septic process.

The mid-20th century marked a pivotal shift towards single-stage resections with primary anastomosis, largely driven by advancements in surgical technique, improved understanding of fluid and electrolyte balance, and the introduction of broad-spectrum antibiotics. Surgeons such as O.H. Wangensteen championed more aggressive, definitive resections. The advent of laparoscopic surgery in the late 20th century, particularly in the 1990s, further revolutionized the field, offering minimally invasive options for elective cases. This led to reduced postoperative pain, shorter hospital stays, and quicker recovery times compared to traditional open approaches. Today, laparoscopic sigmoid colectomy with primary anastomosis is the preferred approach for most elective cases, with robotic assistance also gaining traction.

\section{Clinical Indications}
\begin{itemize}
  \item Complicated diverticulitis, including large or persistent diverticular abscesses (typically >4-5 cm) that fail percutaneous drainage or are inaccessible.
  \item Free perforation with generalized peritonitis (Hinchey III or IV), requiring emergent surgical intervention to control sepsis.
  \item Diverticular fistula formation, such as colovesical (colon to bladder), colovaginal, or coloenteric fistulas, causing debilitating symptoms like pneumaturia or fecaluria.
  \item Bowel obstruction caused by chronic inflammation, stricture, or mass effect from diverticular disease, unresponsive to conservative measures.
  \item Recurrent episodes of acute diverticulitis, particularly after two or more severe attacks, or if attacks are frequent and significantly impact quality of life despite medical management.
  \item Uncontrolled diverticular bleeding that persists despite endoscopic or angiographic interventions, especially if localized to a resectable segment.
  \item Inability to definitively rule out colorectal malignancy, especially when a diverticular mass or stricture mimics cancer on imaging or colonoscopy.
  \item Chronic, smoldering diverticulitis with persistent abdominal pain, altered bowel habits, and low-grade inflammation refractory to medical management.
\end{itemize}

\section*{Clinical Positioning}
Diverticular disease surgery can serve as both a primary and a secondary intervention, depending on the acuity and complexity of the clinical scenario. It is considered a primary surgical treatment for patients presenting with acute, complicated diverticulitis, such as free perforation with generalized peritonitis, large abscesses unresponsive to conservative measures, or acute bowel obstruction. In these emergent situations, surgery is often the first and most definitive treatment to resolve life-threatening complications and achieve source control. For patients with recurrent, uncomplicated diverticulitis or chronic symptoms, surgery is typically a secondary intervention. It is pursued after medical management, including dietary modifications, antibiotics, and lifestyle changes, has failed to control symptoms or prevent further attacks. In these elective cases, surgery is performed to improve the patient's quality of life and prevent future complications, making it a planned, non-urgent procedure.

\section{Surgical Technique \& Procedure}
Diverticular disease surgery is performed under general anesthesia. The patient is typically positioned supine, often with the left arm tucked and the right arm abducted, allowing access for the surgical team. For laparoscopic approaches, the patient may be placed in a steep Trendelenburg position with a left tilt to facilitate exposure of the pelvis and allow small bowel to fall superiorly.

The procedure generally involves the following steps:
\begin{enumerate}
  \item \textbf{Anesthesia and Access:} General anesthesia is induced, and appropriate monitoring is established. For laparoscopic surgery, several small incisions (ports) are made in the abdomen, typically a supraumbilical or umbilical port for the camera, and 2-4 working ports. For open surgery, a midline or left lower quadrant incision is made.
  \item \textbf{Exploration and Mobilization:} The abdomen is explored to assess the extent of diverticular disease, inflammation, and any associated complications (e.g., abscess, fistula). The left colon, including the sigmoid colon, is mobilized from its retroperitoneal attachments. This crucial step involves careful dissection along the white line of Toldt and meticulous identification and protection of the left ureter and gonadal vessels. The splenic flexure may need to be mobilized to achieve a tension-free anastomosis.
  \item \textbf{Identification of Diseased Segment and Margins:} The segment of colon affected by diverticulitis, typically characterized by thickening, inflammation, and diverticula, is identified. The proximal resection margin is usually in healthy, pliable descending colon, free of diverticula, while the distal margin extends into the upper rectum, below the level of the sacral promontory, to ensure removal of all diseased sigmoid colon.
  \item \textbf{Vascular Ligation:} The vascular supply to the diseased segment, primarily the sigmoid branches of the inferior mesenteric artery, is carefully ligated and divided. The IMA itself may be ligated at its origin or distal to the left colic artery, depending on the extent of resection and surgeon preference, always ensuring preservation of the marginal artery of Drummond for collateral flow.
  \item \textbf{Resection:} The diseased segment of the colon is resected using surgical staplers (linear staplers for division) or by hand-sewn techniques. The specimen is then extracted, often through an enlarged port site or the primary incision.
  \item \textbf{Anastomosis or Stoma Creation:} In most elective cases, bowel continuity is restored by creating a primary anastomosis (rejoining the two ends of the colon). This is commonly performed using a circular stapler, creating an end-to-end or end-to-side connection. In emergent cases with significant contamination, severe inflammation, or patient instability, a Hartmann's procedure may be performed, where the distal rectal stump is closed, and the proximal colon is brought out as a diverting colostomy.
  \item \textbf{Leak Test, Irrigation, and Hemostasis:} If an anastomosis is performed, a leak test (e.g., air leak test with saline irrigation) is often conducted. The abdominal cavity is thoroughly irrigated, and meticulous hemostasis is ensured.
  \item \textbf{Closure:} A drain may be placed in some cases, particularly if there was significant contamination, a large abscess cavity, or concern for anastomotic integrity. The incisions are then closed in layers, and port sites are closed to prevent hernia formation.
\end{enumerate}

\section*{Preoperative Workup \& Preparation}
Preparation for diverticular disease surgery is crucial for optimizing outcomes and minimizing risks. For elective cases, a comprehensive preoperative workup is performed. This typically includes a thorough medical history and physical examination, along with routine laboratory tests such as a complete blood count (CBC), electrolyte panel, renal and liver function tests, and coagulation studies. Imaging, usually a computed tomography (CT) scan of the abdomen and pelvis, is essential to confirm the diagnosis, assess the extent of disease, and rule out other pathologies like malignancy. A colonoscopy is often performed preoperatively to evaluate the entire colon and exclude synchronous lesions, unless contraindicated by acute inflammation or obstruction.

Patients will undergo an anesthesia consultation to assess their fitness for general anesthesia and optimize any comorbidities. For elective resections, a mechanical bowel preparation combined with oral antibiotics (e.g., neomycin and metronidazole) is often prescribed to reduce bacterial load in the colon, though its necessity is debated in some centers. Patients are instructed to be NPO (nil per os) for a specified period before surgery, typically after midnight the day before. Intravenous prophylactic antibiotics, usually broad-spectrum, are administered shortly before the surgical incision. Informed consent, detailing the procedure, potential risks, and alternatives, is obtained from the patient. For emergent cases, preparation is expedited, focusing on rapid resuscitation, broad-spectrum antibiotics, and immediate surgical intervention.

\section*{Contraindications \& Precautions}
\textbf{Absolute Contraindications:}
\begin{itemize}
  \item Uncontrolled severe sepsis or septic shock where the patient is too unstable for a definitive resection; in such cases, initial source control (e.g., percutaneous drainage of an abscess or a diverting colostomy without resection) may be necessary to stabilize the patient.
  \item Uncorrectable coagulopathy or severe bleeding diathesis that cannot be managed preoperatively, posing an unacceptable risk of hemorrhage.
  \item Severe, uncompensated end-stage organ failure (e.g., cardiac, pulmonary, renal) where the risks of major abdominal surgery far outweigh any potential benefits, and the patient's life expectancy is severely limited.
\end{itemize}

\textbf{Relative Contraindications \& Precautions:}
\begin{itemize}
  \item Significant comorbidities, such as severe cardiac disease (e.g., recent myocardial infarction, unstable angina), uncontrolled diabetes, or chronic obstructive pulmonary disease (COPD), which significantly increase operative risk. These require careful preoperative optimization and multidisciplinary team involvement.
  \item Extensive intra-abdominal adhesions from previous surgeries, which can make dissection challenging, prolong operative time, and increase the risk of bowel or adjacent organ injury.
  \item Active inflammatory bowel disease (Crohn's disease or ulcerative colitis) coexisting with diverticular disease, as this can complicate diagnosis, surgical planning, and postoperative healing, potentially increasing anastomotic leak risk.
  \item Malnutrition or severe immunosuppression (e.g., from high-dose corticosteroids, chemotherapy, or advanced HIV), which can impair wound healing, increase infection risk, and compromise anastomotic integrity.
  \item Pregnancy, especially in the first and third trimesters, where surgical intervention is generally avoided unless absolutely emergent and life-saving for the mother.
  \item Morbid obesity, which can make laparoscopic access and visualization more challenging, potentially necessitating an open conversion.
\end{itemize}

\section*{Complications \& Risks}
As with any major abdominal surgery, diverticular disease surgery carries inherent risks, which can be broadly categorized:
\begin{itemize}
  \item \textbf{General Surgical Complications:}
    \begin{itemize}
      \item \textbf{Bleeding:} Intraoperative or postoperative hemorrhage, potentially requiring blood transfusion or reoperation.
      \item \textbf{Infection:} Surgical site infection (wound infection), intra-abdominal abscess, or generalized sepsis.
      \item \textbf{Anesthesia Complications:} Adverse reactions to medications, respiratory depression, cardiac arrhythmias, myocardial infarction, or stroke.
      \item \textbf{Deep Vein Thrombosis (DVT) and Pulmonary Embolism (PE):} Blood clots forming in the legs that can travel to the lungs, potentially life-threatening.
      \item \textbf{Ileus:} Prolonged temporary paralysis of the bowel, delaying return of bowel function.
    \end{itemize}
  \item \textbf{Procedure-Specific Risks:}
    \begin{itemize}
      \item \textbf{Anastomotic Leak:} The most serious complication, where the connection between the resected bowel segments fails to heal, leading to leakage of bowel contents into the abdominal cavity. This can cause peritonitis, abscess, or sepsis, often requiring reoperation, diversion with a stoma, or even mortality.
      \item \textbf{Bowel Obstruction:} Can occur early due to ileus or adhesions, or late due to stricture at the anastomosis or recurrent adhesions.
      \item \textbf{Injury to Adjacent Organs:} Accidental damage to the ureters, bladder, small bowel, or major blood vessels (e.g., iliac vessels) during dissection, potentially requiring repair or further surgery.
      \item \textbf{Stoma Complications:} If a colostomy is created, risks include prolapse, retraction, stenosis, parastomal hernia, skin irritation, infection, or high output.
      \item \textbf{Incisional Hernia:} Weakness in the abdominal wall at the incision site, leading to protrusion of abdominal contents, requiring potential repair.
      \item \textbf{Persistent Diverticular Symptoms:} Rarely, symptoms may persist if the entire diseased segment was not removed or if other bowel issues (e.g., irritable bowel syndrome) are present.
      \item \textbf{Fistula Formation:} New fistula formation, though rare, can occur between the bowel and adjacent organs.
      \item \textbf{Sexual Dysfunction:} Particularly in males, due to nerve injury during extensive pelvic dissection, though uncommon with careful technique.
      \item \textbf{Urinary Dysfunction:} Bladder atony or neurogenic bladder due to pelvic nerve injury.
    \end{itemize}
\end{itemize}

\section*{Postoperative Recovery Timeline}
Immediately after diverticular disease surgery, the patient is transferred to the Post-Anesthesia Care Unit (PACU) for close monitoring of vital signs, pain levels, and recovery from anesthesia. Pain management is a priority, often involving intravenous analgesics or patient-controlled analgesia (PCA), transitioning to oral medications as tolerated. Nausea and vomiting are also managed with antiemetics. Early ambulation is strongly encouraged within the first 24 hours to prevent complications such as deep vein thrombosis, pulmonary embolism, and ileus.

Over the first 24-48 hours, the patient's diet is gradually advanced from clear liquids to a soft, low-residue diet as bowel function returns, indicated by the passage of flatus or bowel movements. Intravenous fluids are continued until oral intake is adequate. If a drain was placed, its output is monitored, and it is typically removed when drainage is minimal. For patients with a new stoma, comprehensive stoma care education is initiated by an ostomy nurse. Hospital discharge usually occurs within 3-7 days, depending on the type of surgery (laparoscopic vs. open), the patient's recovery progress, and the absence of complications. Patients are provided with detailed instructions on wound care, pain medication, activity restrictions (e.g., avoiding heavy lifting for 6-8 weeks), and warning signs to watch for. Long-term recovery involves a gradual return to normal activities over several weeks, with full recovery often taking 4-8 weeks. If a temporary stoma was created, a second surgery for stoma reversal is typically planned 3-6 months later, once inflammation has subsided and the patient has fully recovered from the initial procedure.

\section*{Surgical Findings \& Pathology}
During diverticular disease surgery, the surgeon documents detailed intraoperative findings, which are crucial for understanding the extent of the disease and guiding the procedure. These findings typically include the location and extent of diverticula, the degree of inflammation (e.g., pericolic fat inflammation, phlegmon), the presence and size of any abscesses, the location of perforations or fistulas, and the condition of adjacent organs. The operative report will specify the type of procedure performed (e.g., laparoscopic sigmoid colectomy with primary anastomosis, Hartmann's procedure) and any intraoperative complications.

The resected colonic segment is sent for pathological examination. The pathology report is paramount as it confirms the diagnosis of diverticular disease, assesses the severity of inflammation (acute vs. chronic diverticulitis), and, most importantly, definitively rules out the presence of colorectal malignancy, which can sometimes mimic diverticulitis on imaging. A typical pathology report will describe diverticula, often with evidence of chronic inflammation, fibrosis, muscular hypertrophy, and possibly acute inflammatory changes, microperforations, or abscess formation. The absence of malignant cells is a critical finding, providing reassurance to both the patient and the surgical team.

\section{Expected Postoperative Course}
A normal and successful postoperative course following diverticular disease surgery is characterized by an uneventful recovery and the resolution of the symptoms that prompted the surgery. Patients can expect a gradual decrease in incisional pain, which should be well-controlled with oral analgesics within a few days. Bowel function typically returns within 2-5 days, marked by the passage of flatus and then bowel movements, allowing for a progressive advancement of diet. The surgical incision should heal cleanly, without excessive redness, swelling, or purulent drainage. If a stoma was created, it should be healthy, pink, and functioning appropriately, with the patient and caregivers becoming proficient in its management. Patients should experience a significant improvement in their quality of life, with the cessation of recurrent diverticulitis attacks, chronic abdominal pain, or other debilitating complications, enabling them to return to their usual activities without ongoing issues related to their diverticular disease.

\section{Postoperative Care \& Follow-up}
\begin{itemize}
  \item \textbf{Postoperative Clinic Visits:} Typically scheduled 2-4 weeks after discharge to assess wound healing, review the pathology results, discuss the overall recovery, and address any patient concerns.
  \item \textbf{Suture/Staple Removal:} If non-absorbable skin sutures or staples were used, they are usually removed at the first follow-up visit.
  \item \textbf{Stoma Care and Reversal Planning:} For patients with a temporary colostomy, ongoing education on stoma care is provided by an ostomy nurse, and discussions begin regarding the timing and feasibility of stoma reversal surgery, usually 3-6 months post-initial procedure.
  \item \textbf{Colonoscopy:} A full colonoscopy is often recommended 6-12 weeks after surgery (or after stoma reversal) to ensure the entire colon has been adequately screened for polyps or underlying malignancy, especially if a complete preoperative colonoscopy was not possible due to acute inflammation or obstruction.
  \item \textbf{Activity Resumption:} Patients are advised on a gradual return to normal activities, specifically avoiding heavy lifting or strenuous abdominal exercise for 6-8 weeks to allow for complete abdominal wall healing and minimize the risk of incisional hernia.
  \item \textbf{Dietary Advice:} While specific dietary restrictions are often not necessary long-term, some patients may benefit from a high-fiber diet to promote regular bowel habits.
  \item \textbf{Warning Signs:} Patients are advised to contact their surgeon immediately if they experience fever (temperature >101.5$^{\circ}$F or 38.6$^{\circ}$C), worsening or severe abdominal pain, persistent nausea or vomiting, inability to pass gas or stool for more than 24 hours, significant wound drainage, increasing redness or swelling around the incision, or any signs of infection.
\end{itemize}
These comprehensive follow-up measures ensure complete recovery, address any potential long-term issues, and monitor for recurrence or new pathology.

\section*{Epidemiology}
Diverticular disease is highly prevalent in Western populations, with its incidence increasing significantly with age. It is estimated to affect approximately 5\% of individuals under 40, rising to over 50\% in those over 60, and up to 70\% by age 80. While the presence of diverticula (diverticulosis) is common, only about 10-25\% of individuals with diverticulosis will develop diverticulitis, and a smaller proportion (around 15-20\% of those with diverticulitis) will require surgical intervention. The most common indications for surgery are complicated diverticulitis (abscess, perforation, fistula, obstruction) and recurrent episodes that significantly impact quality of life. Emergency surgery for complicated diverticulitis, particularly with free perforation and generalized peritonitis (Hinchey III/IV), carries a significantly higher morbidity and mortality rate (up to 10-20\% mortality) compared to elective surgery (mortality typically less than 1-2\%). There is no significant sex predilection for diverticular disease overall, though some studies suggest a slight male predominance in younger patients and a female predominance in older patients.

\section*{Outcomes \& Prognosis}
The outcomes of diverticular disease surgery are generally favorable, particularly for elective procedures. For elective sigmoid colectomy with primary anastomosis, success rates in resolving symptoms and preventing recurrence are high, often exceeding 90-95\%. Patients typically experience a significant improvement in quality of life, with resolution of chronic pain, recurrent infections, and other debilitating symptoms. Functional outcomes, including return to normal bowel habits, are excellent for most patients. Recurrence of diverticulitis in the remaining colon after a well-performed resection is rare, estimated at 5-10\% over 5-10 years, and often managed medically if it occurs. Prognostic factors influencing outcomes include the patient's age, overall health status, presence of comorbidities, and whether the surgery is performed electively or emergently. Emergency surgery for perforated diverticulitis, especially with generalized peritonitis, carries a higher risk of complications, including anastomotic leak, reoperation, and mortality, but remains a life-saving intervention. For patients undergoing a Hartmann's procedure, the success rate of subsequent stoma reversal is high (70-90\%), though it also carries its own set of risks, including anastomotic leak and wound complications. Landmark studies, such as the APPAC trial, have informed the non-operative management of uncomplicated diverticulitis, but for complicated or recurrent disease, surgery remains the definitive treatment with excellent long-term prognosis.

\section*{References}
\begin{enumerate}
  \item MedlinePlus. Diverticulitis surgery. U.S. National Library of Medicine; 2024. Available from: \url{https://medlineplus.gov/ency/article/007386.htm}
  \item Schwartz SI, Brunicardi FC, Andersen DK, et al. Schwartz's Principles of Surgery. 11th ed. McGraw-Hill Education; 2019.
  \item American Society of Colon and Rectal Surgeons (ASCRS). Clinical Practice Guidelines for the Treatment of Left-Sided Colonic Diverticulitis. Dis Colon Rectum. 2020;63(7):876-887.
  \item Sabiston DC, Townsend CM. Sabiston Textbook of Surgery: The Biological Basis of Modern Surgical Practice. 21st ed. Elsevier; 2021.
\end{enumerate}

\clearpage